\documentclass[aspectratio=169,11pt]{beamer}

\usepackage[T1]{fontenc}
\usepackage[utf8]{inputenc}
\usepackage[french]{babel}
\usepackage{lmodern}
\usepackage{booktabs}
\usepackage{tabularx}
\usepackage{array}
\usepackage{ragged2e}
\usepackage{xcolor}

\usetheme[numbering=fraction,progressbar=frametitle]{metropolis}
\metroset{sectionpage=none,subsectionpage=none}
\setbeamertemplate{navigation symbols}{}

\definecolor{GeneraliRed}{RGB}{172,27,45}
\definecolor{Slate}{RGB}{52,73,94}
\definecolor{SoftGray}{RGB}{245,247,250}

\setbeamercolor{structure}{fg=GeneraliRed}
\setbeamercolor{normal text}{fg=Slate,bg=white}
\setbeamercolor{title}{fg=GeneraliRed}
\setbeamercolor{subtitle}{fg=Slate}
\setbeamercolor{frametitle}{fg=GeneraliRed,bg=white}
\setbeamercolor{block title}{fg=white,bg=GeneraliRed}
\setbeamercolor{block body}{bg=SoftGray}
\setbeamercolor{itemize item}{fg=GeneraliRed}
\setbeamercolor{itemize subitem}{fg=GeneraliRed}

\newcommand{\highlight}[1]{\textcolor{GeneraliRed}{\textbf{#1}}}
\newcommand{\source}[1]{{\scriptsize\textcolor{Slate!65}{[#1]}}}

\title[Thèse CIFRE]{Projet de thèse CIFRE}
\subtitle{Vers un Total Portfolio Approach opérationnel pour Generali France}
\author[Y. Mannai]{Yassine Mannai}
\institute[Generali France]{Generali France \\ Université Paris Dauphine-PSL}
\date[Juin 2026]{Juin 2026}

\begin{document}

\begin{frame}
  \titlepage
\end{frame}

\begin{frame}{Pourquoi ce sujet maintenant ?}
  \begin{itemize}
    \item En 2022, la diversification traditionnelle a cessé de protéger: \highlight{actions et obligations ont baissé ensemble}.
    \item Les dépendances entre facteurs de marché ne sont pas stables: elles changent vite selon le régime macroéconomique.
    \item Pour un assureur, le sujet n'est pas seulement la performance d'actifs, mais le \highlight{pilotage du bilan complet} sous contrainte de capital.
    \item L'enjeu est donc de passer d'une SAA statique par classes d'actifs à une logique plus \highlight{factorielle, dynamique et intégrée}.
    \item L'\highlight{intégration} est ici entendue du point de vue \highlight{actif-passif (ALM)} : il s'agit de piloter l'ensemble du bilan, et non uniquement le portefeuille d'actifs isolément.
  \end{itemize}
\end{frame}

\begin{frame}[shrink=15]{Constat sur le processus actuel}
  \begin{columns}[T,totalwidth=\textwidth]
    \column{0.48\textwidth}
    \begin{block}{Chaîne de décision aujourd'hui}
      \begin{itemize}
        \item GSE avec corrélations historiques moyennes
        \item SCR lourd via modèles actuariels et proxy LSMC
        \item Optimisation SAA revue à fréquence annuelle
        \item Mise en œuvre opérationnelle partiellement manuelle
      \end{itemize}
    \end{block}

    \column{0.48\textwidth}
    \begin{block}{Limites pour l'équipe SAA}
      \begin{itemize}
        \item Faible prise en compte des ruptures de régime
        \item Temps de calcul incompatible avec des boucles rapides
        \item Allocation et couverture pensées séparément
        \item Actifs privés lissés : risque économique sous-estimé
        \item Scénarios de passif \highlight{statiques} (pas de stress dynamique \emph{best/worst case})
        \item Boucles SAA~$\leftrightarrow$~Risques (SCR stock, optimisation, recalcul) à unifier dans un \highlight{optimiseur unique}
      \end{itemize}
    \end{block}
  \end{columns}
\end{frame}

\begin{frame}[shrink=10]{Ambition de la thèse}
  \begin{block}{Question centrale}
    Comment rendre le \highlight{Total Portfolio Approach} exploitable pour un assureur sous Solvabilité II, en tenant compte des changements de régime, du coût en capital et des besoins de couverture ?
  \end{block}

  \begin{itemize}
    \item Construire une lecture commune des risques via des \highlight{facteurs économiques partagés} entre actifs et passifs.
    \item Répondre vite aux questions \highlight{``que faire si le régime change ?''}.
    \item Livrer des \highlight{modules industrialisés intégrés à GPS} (\emph{Group Portfolio Solutions}, l'outil SAA de Generali).
    \item \highlight{Fil rouge mathématique} : contrôler la \highlight{propagation d'erreur} le long de la chaîne \emph{filtre $\to$ générateur $\to$ proxy $\to$ contrôle} — c'est ce qui unifie les trois axes en une seule thèse, et non trois projets juxtaposés.
  \end{itemize}

  {\footnotesize Fonction d'utilité \highlight{multi-dimensionnelle} : rendement, \emph{carry}, matching ALM, liquidité, SCR. Le TPA permet d'\highlight{élargir cette fonction d'utilité} au-delà du seul rendement — là où les fonds souverains (Future Fund, CPP) l'appliquent au rendement total, Generali en adapte la logique au \highlight{bilan assurantiel}. Cette utilité multi-objectifs se formalise comme une \highlight{mesure de risque shortfall multivariée}, réductible à un problème unidimensionnel. \source{Doldi, Frittelli, Rosazza Gianin 2023}}
\end{frame}

\begin{frame}[shrink=15]{Vision cible pour Generali France}
  \begin{columns}[T,totalwidth=\textwidth]
    \column{0.5\textwidth}
    \begin{block}{Les facteurs TPA à piloter}
      \begin{itemize}
        \item Duration (taux)
        \item Inflation
        \item Croissance
        \item Liquidité / carry
        \item \highlight{Facteur transverse actif-passif} : la courbe des taux reprice obligations \emph{et} BEL
      \end{itemize}
    \end{block}

    \column{0.5\textwidth}
    \begin{block}{Ce que cela permet}
      \begin{itemize}
        \item Mesurer l'impact marginal d'une décision sur le bilan total
        \item Comparer actifs cotés, actifs privés et couvertures
        \item Préparer des plans d'action avant les phases de stress
        \item Générer un \highlight{surplus de rendement sous contrainte de bilan} — effet secondaire, non objectif principal
      \end{itemize}
    \end{block}
  \end{columns}

  \vspace{0.1cm}
  \centering
  {\footnotesize \highlight{Idée directrice :} étendre le TPA au passif et piloter le bilan par facteurs sous contrainte de capital, plutôt que par silo d'actifs.}
\end{frame}

\begin{frame}[shrink=20]{Axe 1: génération conditionnelle au régime, avec garanties}
  \begin{itemize}
    \item Identifier en temps réel le \highlight{régime macro-financier} via un filtre séquentiel (SMC).
    \item Simuler les facteurs TPA conditionnellement au régime, avec dépendances de queue non linéaires: d'abord un modèle génératif flexible (GAN), puis un modèle à \highlight{garanties de convergence} fondé sur le transport optimal entropique — \highlight{ponts de Schrödinger à sauts} capturant queues épaisses, sauts et changements de régime que les modèles purement diffusifs manquent. \source{De Marco, Pham, Zanni 2026}
    \item Valider les dépendances par \highlight{tests de copules} $\Rightarrow$ auditabilité réglementaire.
    \item Lire les dépendances entre facteurs et entités comme un \highlight{réseau} (contagion, mesures de type DebtRank) pour le risque systémique en environnement dégradé. \source{Bardoscia et al.\ 2021}
    \item Traiter le \highlight{retard de valorisation des actifs privés} (\emph{valuation lag}) et la \highlight{stochasticité réelle des trajectoires économiques}, masquée par le lissage comptable — l'enjeu est de réconcilier perception comptable et réalité économique du risque.
  \end{itemize}

  \vspace{0.1cm}
  {\footnotesize \highlight{Contribution théorique visée :} bornes de convergence ($W_2$) de la génération conditionnelle lorsque le régime est \emph{estimé} — propager l'erreur d'inférence du filtre dans la borne, cas ouvert dans la littérature.}

  \begin{block}{Livrable métier (module GPS \textemdash{} \textit{factor engine})}
    Un moteur de scénarios réalistes pour la SAA et le TPA, plus un \highlight{scénario central} (médoïde) pour les comités et le budget, intégré nativement dans \highlight{GPS}.
  \end{block}
\end{frame}

\begin{frame}[shrink=15]{Axe 2: proxy SCR rapide, avec bornes d'erreur contrôlées}
  \begin{itemize}
    \item Apprendre un proxy supervisé du \highlight{SCR} (réseau de neurones) à partir des facteurs, du régime et des états passif, dont le lag des actifs privés. Référence et validation par \highlight{simulation imbriquée / Monte Carlo multi-niveaux}. \source{Bourgey 2020}
    \item \highlight{Pondération adaptative}: la précision est concentrée là où elle compte, près du seuil réglementaire et en régime de stress, avec des bornes d'erreur par régime.
    \item Passer de \highlight{plusieurs heures à moins de 5\,ms} par évaluation.
    \item Intégrer le proxy comme \highlight{contrainte de capital explicite} dans l'optimisation, pas seulement comme objet de reporting.
  \end{itemize}

  {\footnotesize \highlight{Contribution théorique visée :} analyse d'erreur d'un estimateur imbriqué \emph{avec biais de lissage} (lag des actifs privés), qui sort des hypothèses standard du multi-niveaux — bornes de généralisation par régime.}

  \begin{block}{Livrable métier (module GPS \textemdash{} \textit{SCR proxy})}
    Un SCR quasi instantané, utilisable dans des boucles de décision, d'allocation et de couverture, exposé comme service au sein de \highlight{GPS}.
  \end{block}
\end{frame}

\begin{frame}[shrink=10]{Axe 3: politique jointe allocation–couverture (contrôle stochastique)}
  \begin{itemize}
    \item Formuler la décision comme un \highlight{problème de contrôle stochastique} sous \highlight{contrainte de risque dynamique} (mesures time-consistent, caractérisation \highlight{BSDE}, séparation en fonds) ; le \highlight{deep hedging} sert de \emph{schéma numérique} d'approximation, pas de boîte noire. \source{Moreno-Bromberg, Pirvu, Réveillac 2013}
    \item \highlight{Identifier automatiquement l'objectif dominant} par portefeuille via la \highlight{contribution marginale au capital} (allocation d'Euler/gradient, cas convexe non différentiable). \source{Centrone, Rosazza Gianin 2025}
    \item Arbitrer sous contraintes de capital (proxy SCR), de liquidité et d'engagements irrévocables des actifs privés, pour réduire le \highlight{risque de queue} du bilan (CVaR du déficit).
  \end{itemize}

  {\footnotesize \highlight{Contribution théorique visée :} caractérisation de la politique jointe quand la contrainte est un \emph{proxy estimé} (axe 2) et la dépendance \emph{à changement de régime} (axe 1), avec actifs irrévocables — hors du cadre de Réveillac (2013).}

  \begin{block}{Livrable métier (module GPS \textemdash{} \textit{allocator})}
    Une \highlight{fonction de réaction d'allocation}: un plan what-if calculé ex ante, avec déclenchement explicite, intégré à \highlight{GPS} et connecté à l'interface d'exécution GenAM.
  \end{block}
\end{frame}

\begin{frame}{Risque de marché et risque actuariel}
  \begin{itemize}
    \item \highlight{Une meilleure lecture des risques de dépendance} en environnement dégradé.
    \item \highlight{Des arbitrages plus rapides} entre rendement, SCR et coût de couverture.
    \item \highlight{Un cadre commun} entre investissement, risque et actuariel autour de facteurs de bilan.
    \item \highlight{Un outil industrialisable} pour préparer des décisions et challenger des allocations.
  \end{itemize}
\end{frame}

\begin{frame}{Feuille de route sur 3 ans}
  \renewcommand{\arraystretch}{1.25}
  \begin{tabularx}{\textwidth}{>{\RaggedRight\arraybackslash}p{1.7cm}>{\RaggedRight\arraybackslash}X}
    \toprule
    \textbf{Année} & \textbf{Objectif principal et livrable} \\
    \midrule
    Année 1 & Base de données facteurs (20 ans, lag actifs privés), détection de régimes en ligne, GAN conditionnel validé par copules, scénario central. \highlight{Module \textit{factor engine}} intégré à GPS. \\
    Année 2 & Proxy SCR conditionné au régime (< 5\,ms), bornes d'erreur, intégration comme contrainte d'optimisation; SGM à garanties $W_2$. \highlight{Module \textit{SCR proxy}} exposé dans GPS. \\
    Année 3 & Politique jointe allocation / couverture, backtests multi-régimes (2008, 2020, 2022), interface GenAM. \highlight{Module \textit{allocator}} intégré à GPS. \\
    \bottomrule
  \end{tabularx}
\end{frame}

\begin{frame}{}
  \vfill
  \centering
  \Large \highlight{Merci !}
  \vfill
\end{frame}

\end{document}